%%%%%%%%%%  scan idx 276  |  folio 265  %%%%%%%%%%

\subsection*{4. Extensiones y biextensiones de esquemas en grupos
lisos y conexos sobre un cuerpo}

Recordemos el lema siguiente, debido a Rosenlicht [3 bis, Proposición
3].

\medskip
\noindent\underline{Lema 4.1.} \underline{Sean $k$ un cuerpo, y $X$
e $Y$ dos $k$-esquemas separables, geométricamente conexos y de tipo
finito, provistos respectivamente de puntos $k$-racionales $e_X$ y
$e_Y$. Entonces todo morfismo}
\[
 X\times Y\longrightarrow\underline G_m
\]
\underline{que es trivial sobre los dos subesquemas $X\times e_Y$ y
$e_X\times Y$ es trivial.}

Utilizaremos este resultado únicamente cuando $X$ e $Y$ sean esquemas
en grupos sobre $k$; en tal caso, las hipótesis implican que $X$ e $Y$
son de hecho lisos sobre $k$. Por inducción, 4.1 implica también el
resultado análogo para un número cualquiera de factores. Podemos, por
tanto, aplicar los resultados de VII 1.3 y VII 2.9 y obtener las dos
proposiciones siguientes.

\medskip
\noindent\underline{Proposición 4.2.} \underline{Sea $P$ un grupo
liso y conexo sobre el cuerpo $k$. El funtor natural de la categoría
de las extensiones de grupos de $P$ por $\underline G_m$ a la categoría
de los módulos invertibles sobre $P$ rigidificados en el elemento
neutro $e_P$ es plenamente fiel}. \underline{Su imagen esencial está formada por
los módulos invertibles rigidificados $\mathbb L$ sobre $P$ para los
que}
\[
 \pi^*\sisi{P}{\mathbb L}\simeq
 \pr_1^*\sisi{P}{\mathbb L}\sisi{}{\otimes}
 \pr_2^*\sisi{P}{\mathbb L},
\]
\underline{donde $\pi:P\times P\longrightarrow P$ es la ley de grupo
de $P$, y $\pr_1,\pr_2$ son las dos proyecciones.}

\medskip
\noindent\underline{Proposición 4.3.} \underline{Sean $P$ y $Q$ dos
grupos algebraicos lisos y conexos sobre el cuerpo $k$. El funtor
natural de la categoría de las biextensiones de $(P,Q)$ por
$\underline G_m$ a la categoría de los módulos invertibles sobre
$P\times Q$ birigidificados con respecto a las secciones identidad
$e_P$ y $e_Q$ de $P$ y $Q$ es plenamente fiel; además, ambas
categorías son rígidas.}

%%%%%%%%%%  scan idx 277  |  folio 266  %%%%%%%%%%

\underline{Un módulo invertible birigidificado $\mathbb L$ sobre
$P\times Q$ procede de una biextensión si y solo si sus imágenes
inversas sobre $P\times P\times Q$ y $P\times Q\times Q$ por
$\pi_P\times\id_Q$ y $\id_P\times\pi_Q$ satisfacen el ``teorema del
cubo'' ([1] o [2]).}

Esta última condición no es sino la traducción de las condiciones de
VII 2.9.3 b), en la forma explícita dada en VII 3.5 b).

\medskip
\noindent 4.4. No sé si la condición del cubo de 4.3 se satisface
automáticamente. Al menos se sabe [2, IV 2.6] que, para un
$\mathbb L$ dado, siempre existe un entero $n\geq1$, potencia del
exponente característico de $k$, tal que $\mathbb L^{\otimes n}$
satisface esa condición y, por tanto, procede de una biextensión de
$(P,Q)$ por $\underline G_m$. También es siempre posible encontrar
una extensión puramente inseparable $k'$ de $k$ tal que la condición
requerida se haga verdadera después de la extensión de base
$k\longrightarrow k'$; tras esta extensión, por tanto, $\mathbb L$ es
un módulo invertible birigidificado procedente de una biextensión de
$(P,Q)$ por $\underline G_m$.

En particular, cuando $k$ es perfecto, la condición de 4.3 se
satisface siempre. Este último hecho, que implica formalmente las
afirmaciones precedentes mediante un argumento de traza, también puede
verse aplicando 4.7 más abajo; cf. 4.9.
